% Light styling for the pandoc/xelatex notes in docs/.
\usepackage[table]{xcolor}
\usepackage{colortbl}
\definecolor{accent}{HTML}{0072B2}
\definecolor{ink}{HTML}{263238}
\definecolor{muted}{HTML}{657278}
\definecolor{panel}{HTML}{F2F6F8}

\color{ink}
\usepackage{microtype}
\setlength{\parskip}{0.55em}
\setlength{\parindent}{0pt}
\linespread{1.08}

% Headings: sans, accent colour, modest sizes.
\usepackage{titlesec}
\titleformat{\section}{\Large\bfseries\color{accent}}{}{0pt}{}[{\color{accent!35}\titlerule[0.6pt]}]
\titleformat{\subsection}{\large\bfseries\color{ink}}{}{0pt}{}
\titleformat{\subsubsection}{\normalsize\bfseries\color{muted}}{}{0pt}{}
\titlespacing*{\section}{0pt}{1.6em}{0.7em}
\titlespacing*{\subsection}{0pt}{1.2em}{0.4em}
\titlespacing*{\subsubsection}{0pt}{0.9em}{0.2em}

% Title block.
\makeatletter
\renewcommand{\maketitle}{%
  \begin{minipage}{\linewidth}
    {\color{accent}\rule{2.2cm}{3pt}}\par\vspace{0.6em}
    {\fontsize{22}{26}\selectfont\bfseries \@title\par}\vspace{0.35em}
    {\large\color{muted} \notesubtitle\par}\vspace{0.6em}
    {\small\color{muted} \@author\quad·\quad \@date\par}
  \end{minipage}\par\vspace{1.2em}}
\makeatother

% Running header and footer.
\usepackage{fancyhdr}
\pagestyle{fancy}
\fancyhf{}
\renewcommand{\headrulewidth}{0pt}
\fancyhead[L]{\footnotesize\color{muted} Silicon sampling au Québec · premiers résultats}
\fancyhead[R]{\footnotesize\color{muted} \thepage}
\fancypagestyle{plain}{\fancyhf{}\renewcommand{\headrulewidth}{0pt}}

% Tables: more air, muted rules.
\renewcommand{\arraystretch}{1.25}
\usepackage{booktabs}
\setlength{\aboverulesep}{0.3ex}
\setlength{\belowrulesep}{0.5ex}
\arrayrulecolor{muted!60}

% Lists.
\usepackage{enumitem}
\setlist[itemize]{leftmargin=1.3em, itemsep=0.25em, topsep=0.3em, label={\color{accent}\textbullet}}

% Boxes used through raw LaTeX in the markdown.
\usepackage[most]{tcolorbox}
\newtcolorbox{encadre}{colback=panel, colframe=panel, boxrule=0pt,
  leftrule=3pt, colframe=accent, sharp corners, left=10pt, right=10pt, top=6pt, bottom=6pt,
  before skip=0.8em, after skip=0.8em}

\usepackage{caption}
\captionsetup{font={small,color=muted}, labelformat=empty}
\hypersetup{colorlinks=true, linkcolor=accent, urlcolor=accent}

% Prompt examples: wrapped, lightly shaded code blocks (```text in the markdown).
\usepackage{fvextra}
\RecustomVerbatimEnvironment{Highlighting}{Verbatim}{breaklines,breakanywhere,fontsize=\footnotesize,commandchars=\\\{\}}
\definecolor{shadecolor}{HTML}{F2F6F8}

% Table highlighting (\cellcolor in the markdown tables).
\definecolor{best}{HTML}{CDEBDD}
\definecolor{worse}{HTML}{F8D7D2}
\definecolor{ref}{HTML}{ECEFF1}
